\begin{description}
\item[(a)] From Wien's law, $b=T\lambda$, where $b$ is Wien's constant, $T$ is the surface temperature, and $\lambda$ is the wavelength corresponding to the maximum peak of radiation. Hence
\[
b=T_A\lambda_A=T_B\lambda_B\quad\Rightarrow\quad\frac{T_A}{T_B}=\frac{\lambda_B}{\lambda_A}.
\]\hfill(2 points)
The Stefan--Boltzmann formula gives $L=4\pi R^2\sigma T^4$, so
\[
\frac{L_A}{L_B}=\left(\frac{R_A}{R_B}\right)^2\left(\frac{T_A}{T_B}\right)^4
=\left(\frac{R_A}{R_B}\right)^2\left(\frac{\lambda_B}{\lambda_A}\right)^4
=2^2\left(\frac{400}{650}\right)^4\approx0.57.
\]\hfill(2 points)
The luminosity--absolute magnitude relationship is
\[
\frac{L_A}{L_B}=10^{0.4(M_B-M_A)},\qquad
\Delta M=M_A-M_B=-\frac{1}{0.4}\log_{10}\frac{L_A}{L_B}
=-\frac{1}{0.4}\log_{10}0.57\approx0.60.
\]\hfill(2 points)
Also accept 0.61, possible due to rounding.
\item[(b)] Denote by $L=L_A+L_B$ the total luminosity and by $M=0$ the absolute magnitude of the system. Comparing with the Sun,
\[
\frac{L}{L_\odot}=10^{0.4(M_\odot-M)}
=\frac{L_A}{L_\odot}+\frac{L_B}{L_\odot}
=10^{0.4(M_\odot-M_A)}+10^{0.4(M_\odot-M_B)}.
\]\hfill(1 point)
Thus
\begin{align*}
10^{-0.4M}&=10^{-0.4M_A}+10^{-0.4M_B}\\
&=10^{-0.4(M_B+\Delta M)}+10^{-0.4M_B}.
\end{align*}\hfill(1 point)
It follows that
\begin{align*}
10^{-0.4M}&=10^{-0.4M_B}(10^{-0.4\Delta M}+1),\\
-0.4M&=-0.4M_B+\log_{10}(10^{-0.4\Delta M}+1),\\
M_B&=M+\frac{1}{0.4}\log_{10}(10^{-0.4\Delta M}+1)\\
&=0+\frac{1}{0.4}\log_{10}(10^{-0.4\cdot0.61}+1)\approx0.49.
\end{align*}\hfill(1 point)
Then $M_A=M_B+\Delta M=0.49+0.61=1.10$. \hfill(1 point)

\textbf{Alternative solution to (b).} With the same notation,
\[
M-M_A=-2.5\log\frac{L_A+L_B}{L_A}=-2.5\log\left(1+\frac{L_B}{L_A}\right).
\]\hfill(2 points)
Hence $0-M_A=-2.5\log(1+1/0.57)$ and $M_A=2.5\log(2.74)=1.10$. \hfill(1 point)
Using the previously calculated difference, $M_B=M_A-\Delta M=1.1-0.60=0.50$. \hfill(1 point)
\end{description}
